\documentclass[10pt]{article}
\usepackage{amsmath}
\usepackage{booktabs}
\usepackage{multicol}
\usepackage{enumitem}
\usepackage{array}
\usepackage[landscape,margin=0.30in,top=0.32in]{geometry}
\usepackage{microtype}
\usepackage{titlesec}

\setlength{\columnseprule}{0.3pt}
\setlength{\columnsep}{12pt}
\setlength{\parskip}{0.5pt}
\setlength{\parindent}{0pt}
\setlength{\abovedisplayskip}{2pt}
\setlength{\belowdisplayskip}{2pt}
\setlength{\abovedisplayshortskip}{0pt}
\setlength{\belowdisplayshortskip}{0pt}

\titlespacing*{\section}{0pt}{0pt}{0pt}

\newcommand{\sect}[1]{%
  \vspace{4pt}\noindent{\footnotesize\bfseries #1}%
  \vspace{0.5pt}\\\rule{\columnwidth}{0.3pt}\vspace{1pt}}
\newcommand{\con}[1]{\noindent{\scriptsize\textit{$\triangleright$ #1}}\\}

\begin{document}
\begin{center}
{\small\bfseries 15.535 BAUFS Formula \& Concept Sheet --- Sessions 1--11 \hfill MIT Sloan, Spring 2026}
\end{center}
\vspace{-5pt}\rule{\linewidth}{0.4pt}\vspace{2pt}

\begin{multicols}{3}
\footnotesize

%----- STRATEGY -----
\sect{Strategy Framework (S1)}
\textbf{Porter's Five Forces}: Rivalry (many competitors, low diff., high fixed costs); Entrants (capital, scale, patents, regulation); Substitutes \& switching costs; Buyer power (concentrated, undiff.); Supplier power (few, high switching costs). Strong forces $\Rightarrow$ lower margins.\\
\textbf{Economic Attributes}: Demand (price sensitivity, cyclicality, seasonality); Supply (diff., barriers); Production (capital vs.\ labor intensity); Marketing (B2B/B2C, push/pull); Assets (short vs.\ long-term).\\
\con{Strategy sets the benchmark for expected ratios. Deviations require explanation. Key Q: Is performance sustainable or transitory?}

%----- RATIOS / DUPONT -----
\sect{Ratios \& DuPont (S2)}
\textbf{ROE} $=$ NI/Avg.Eq \quad \textbf{ROA} $=$ NI/Avg.TA\\
\textbf{DuPont:} $\text{ROE} = \frac{NI}{Sales}\times\frac{Sales}{\text{Avg.TA}}\times\frac{\text{Avg.TA}}{\text{Avg.Eq}}$ = Margin $\times$ Turnover $\times$ Leverage\\
Profitability: Gross/EBITDA/EBIT/Net margins\\
Activity: AR/Inv/Payables/Fixed Asset Turnover\\
Leverage: D/TA, D/E \quad Liquidity: Current, Quick, Int.Cov.\\
Valuation: P/E, P/S, P/B, TEV/EBITDA\\
\textbf{NCI rule}: use \textit{total} NI+Equity \textit{or} NI-to-common+Common Eq. Never mix.\\
\textbf{Component isolation}: assign one co.'s component to other's formula to quantify impact (``proforma recast'').\\
\con{Asset turnover has more permanent effect on ROA than margin (Soliman 2008).}

%----- PIOTROSKI -----
\sect{Piotroski F-Score (S2) --- 9 binary signals}
\begin{enumerate}[noitemsep,topsep=0pt,leftmargin=*,label=\arabic*.,itemsep=0pt]
\item ROA $> 0$ \quad 2. CFO $> 0$ \quad 3. CFO $>$ NI
\item $\Delta$ROA $> 0$ \quad 5. $\Delta$Gross Margin $> 0$
\item $\Delta$Asset Turnover $> 0$
\item $\Delta$(LTD/TA) $< 0$ (incl.\ curr.\ portion)
\item $\Delta$Current Ratio $> 0$ \quad 9. No new equity to public
\end{enumerate}
\textbf{7--9} strong; \textbf{0--3} weak. ROA = NI (bef.\ disc.ops) / Beg.TA\\
\con{Designed for low P/B neglected stocks. Works in portfolio (idiosyncratic shocks cancel); not a commodity/macro predictor. High F $\Rightarrow$ profitability increases, positive earnings surprises, stock outperformance.}

%----- ACCRUALS -----
\sect{Accruals \& Earnings Quality (S4)}
Total Accruals $= NI - CFO$\\
Accrual Ratio (CFA): $(NI-CFO-CFI)/\text{Avg.TA}$; also $(NI-CFO)/\text{Avg.TA}$\\
WC accruals: $\Delta$CA $-$ $\Delta$CL (ex. cash/ST debt); volatile, short-term\\
Non-WC: D\&A, def.\ taxes; stable, larger\\[1pt]
\textbf{Sign convention}: CFS \textit{backs out} accruals from NI $\Rightarrow$ accrual $=$ \textit{opposite sign} of CFS line.\\
\textit{e.g.} Inv.\ increase: $(-98)$ on CFS $\Rightarrow$ accrual $= +98$.\\[1pt]
\textbf{Non-cash test}: IS charge $=$ CFS add-back $\Rightarrow$ 100\% non-cash that period.\\[1pt]
\textbf{Inventory accrual $\neq$ $\Delta$BS Inventory} (acquisitions, write-offs, FX can differ).\\[1pt]
\textbf{CF manipulation}: Tyco (investing $\to$ operating); Enron (investing $\to$ operating); WorldCom (operating $\to$ investing); Parmalat (falsified accounts).\\
\con{Positive accruals $\Rightarrow$ lower quality. Accr. reverse over time. Managers shift timing, not create profits.}

%----- DAYS -----
\sect{Days Metrics (S4--5)}
Days Inv $= 365\div(\text{COGS}/\text{Avg.Inv})$ \quad Days AR $= 365\div(S/\text{Avg.AR})$\\
Days AP $= 365\div(\text{COGS}/\text{Avg.AP})$ \quad Days Unearn $= 365\div(S/\text{Avg.Unearn})$\\
\con{Declining Days AR + booming sales $=$ healthy. Rising Days Inv + backlog $=$ ramp-up (not necessarily bad). Stable Days Unearn $=$ mix stable; new high-days products offsetting low-days ones.}

%----- COOKIE JAR -----
\sect{Cookie Jar Reserves (S5--7)}
Doubtful\%: Allow/$\text{AR}_g$ \; ($\text{AR}_g=\text{AR}_n+$Allow)\\
Returns/Warranty\%: Addtn to Reserve/Sales\\
Dep.Life $=$ Avg.Gross PPE/Dep.Exp \; (Gross $=$ Net $+$ Accum.); Rem.life $=$ Net/Dep.Exp\\
LIFO$\to$FIFO: $\Delta$LIFO Res $=$ COGS$_L -$ COGS$_F$; FIFO COGS $=$ LIFO COGS $-\,\Delta$LIFO Res\\
LIFO layer liquidation: drawing down inv.\ levels forces old low-cost inventory into COGS $\Rightarrow$ artificially \textit{reduces} COGS $\Rightarrow$ inflates profits. Firms must disclose this effect.\\
DTA Allowance: required when DTA ``more likely than not'' not realizable; $\uparrow$ allowance $\Rightarrow$ $\uparrow$ tax expense.\\
\con{Aggressive accruals foreshadow reversals. Accounting diffs can be legitimate; aggressive on one can be conservative on another.}

%----- SPECIAL ITEMS -----
\sect{Special Items \& Non-GAAP (S6--7)}
\textbf{Core/recurring}: SG\&A, R\&D, royalties, amortization (same sign, every year).\\
\textbf{Special}: Restructuring, impairments, M\&A costs, litigation, asset G/L, disc.ops.\\
\textbf{Gray areas}: Litigation (winding down vs.\ ongoing); Contingent consideration (recurring if M\&A is core strategy, special if one-off). Justify with pattern + strategy.\\
\textbf{Special items on CFS too}: Debt extinguishment, investment G/L.\\
\textbf{Adj. rule}: Back out from NI; adjust tax; ignore BS impact. ``Back out'' $\neq$ ``ignore.''\\
\textbf{Non-GAAP}: Street earnings; excludes specials $+$ sometimes recurring (stock comp). You decide if reasonable.\\
\textbf{OCI} (bypasses IS $\to$ equity): AFS sec., hedges, pension G/L, FX translation, IFRS reval.\\
\con{Negative specials $=$ cockroaches: more will follow. Non-recurring $\neq$ short-lived. Turnaround: specials $\Rightarrow$ corrective actions $\Rightarrow$ improvement eventually.}

%----- Z-SCORE -----
\sect{Bankruptcy \& Z-Scores (S3)}
\textbf{Z} (public mfg): $1.2\frac{WC}{TA}+1.4\frac{RE}{TA}+3.3\frac{EBIT}{TA}+0.6\frac{MVE}{BVL}+1.0\frac{S}{TA}$\\
$Z>3.0$ safe; $1.8$--$3.0$ grey; $Z<1.8$ distress\\[2pt]
\textbf{Z$''$} (any non-fin): $3.25+6.56\frac{WC}{TA}+3.26\frac{RE}{TA}+6.72\frac{EBIT}{TA}+1.05\frac{BVE}{BVL}$\\
$Z''>5.85$ safe; $4.35$--$5.85$ grey; $Z''<4.35$ distress. Maps to bond ratings.\\
EBIT $=$ pretax income $+$ int.exp. \; MVE $=$ price $\times$ shares at FYE.\\EBIT/TA and S/TA use \textit{year-end} TA (not avg); other Z vars unspecified. \textit{Z article cutoffs wrong; use slide values.}\\[2pt]
{\scriptsize\begin{tabular}{@{}l@{\ }c@{\ }c@{\ }c@{\ }c@{}}
\toprule
Event & WC/TA & RE/TA & EBIT/TA & BVE/BVL\\
\midrule
IPO/equity raise & $\uparrow$ & $\uparrow^*$ & $\uparrow^*$ & $\uparrow$\\
Net loss & $\circ$ & $\downarrow$ & $\downarrow$ & $\downarrow$\\
New debt & ${\uparrow}$ST & $\circ$ & $\circ$ & $\downarrow$\\
Asset write-off & $\circ$ & $\downarrow$ & $\downarrow$ & $\downarrow$\\
\bottomrule
\end{tabular}}\\[1pt]
\con{IPO: CA$\uparrow$, TA$\uparrow$, BVE$\uparrow$ (via APIC); EBIT and RE numerators unaffected. $^*$For startup (RE$<$0, EBIT$<$0): ratios increase toward zero. For profitable co (RE$>$0, EBIT$>$0): RE/TA$\downarrow$, EBIT/TA$\downarrow$. Cash cushion $\neq$ fundamental health.}

%----- M-SCORE -----
\sect{Beneish M-Score (S8)}
{\scriptsize
$M=-4.84+0.92\,DSRI+0.528\,GMI+0.404\,AQI+0.892\,SGI$\\
$\phantom{M=}+0.115\,DEPI-0.172\,SGAI-0.327\,LVGI+4.679\,TATA$
}\\[1pt]
{\scriptsize\begin{tabular}{@{}l@{\ }l@{}}
\toprule
Var & Formula (num yr$_t$ / denom yr$_{t-1}$ unless noted)\\
\midrule
DSRI & $(AR/S)_t/(AR/S)_{t-1}$\\
GMI & $GM_{t-1}/GM_t$\\
AQI & $(1{-}(CA{+}PPE_n)/TA)_t/(\cdots)_{t-1}$\\
SGI & $S_t/S_{t-1}$\\
DEPI & $(D/(D{+}PPE_n))_{t-1}/(\cdots)_t$\\
SGAI & $(SGA/S)_t/(SGA/S)_{t-1}$\\
LVGI & $((LTD{+}CL)/TA)_t/(\cdots)_{t-1}$\\
TATA & $(NI_\text{bef.disc}{-}CFO)/TA_t$\\
\bottomrule
\end{tabular}}\\[1pt]
Higher M $=$ more risk. 2+ vars near manipulator col $=$ red flag.\\
Set AQI$=1$ if $(CA{+}PPE)=TA$ in $t{-}1$; PPE incl.\ capit.\ leases.\\[1pt]
\textbf{Common errors}: AQI: use net PPE (not bldgs); CA$+$PPE / TA (not CL). DEPI: depreciation only (not D\&A). LVGI: LTD $+$ CL in numerator / TA. TATA: divide by TA.\\[1pt]
\textbf{Why errors cancel}: ratio-of-ratios $\Rightarrow$ num/denom errors offset; TATA small $\Rightarrow$ big rel.\ error, small abs.\ effect; LVGI negative coeff.\ moderates.\\
\con{Intersegment transactions net to zero in consol.\ FS $\Rightarrow$ M-Score cannot detect segment-level manipulation.}

%----- DECHOW -----
\sect{Dechow F-Score (S8)}
{\scriptsize
$\hat{y}=-7.893+0.790r+2.518c_r+1.191c_i+1.979s$\\
$\phantom{\hat{y}=}+0.171c_s-0.932c_\rho+1.029\,\text{iss}$
}\\
Prob $= e^{\hat{y}}/(1+e^{\hat{y}})$; \; F $=$ Prob$/0.0037$\\

F$>1$ above normal; F$>1.85$ substantial; F$>2.45$ high\\[2pt]
{\scriptsize\begin{tabular}{@{}l@{\ }l@{}}
\toprule
Var & Definition\\
\midrule
$r$ (rsst\_acc) & $\Delta$Non-cash NOA/Avg.TA\\
$c_r$ (ch\_rec) & $\Delta AR_n$/Avg.TA\\
$c_i$ (ch\_inv) & $\Delta$Inv/Avg.TA\\
$s$ (soft\_assets) & $(TA{-}PPE_n{-}Cash)/TA$\\
$c_s$ (ch\_cs) & $\%\Delta(S{-}\Delta AR)$\\
$c_\rho$ (ch\_roa) & $\Delta(NI/\text{Avg.TA})$\\
iss (issue) & $=1$ if LTD or equity issued\\
\bottomrule
\end{tabular}}\\[1pt]
Non-cash NOA $=$ Equity $-$ Pref.Stk $-$ Cash. NI bef.\ disc.ops; PPE incl.\ cap.\ leases; equity incl.\ SBC.\\
\con{Score mainly driven by soft\_assets and issue. Asset-light (low D\&A) $\Rightarrow$ CFO $<$ NI is \textit{not} automatically a bad signal (D\&A is usually the biggest accrual).}

%----- BENFORD -----
\sect{Benford's Law (S8)}
Distribution of the occurrence of 1st digit of numbers in a group of numbers. Tends towards lower numbers. \\
$P(d) = \log_{10}(1+1/d)$\\[2pt]
Applies to naturally occurring data (revenues, populations). Does NOT apply to assigned numbers (zip codes, phone numbers). Deviations $=$ manipulation signal.\\
\con{Fraudsters can't generate convincingly random-looking numbers.}

%----- VALUATION -----
\sect{Valuation Multiples (S2)}
Trailing P/E $= P/$TTM EPS; \; Fwd P/E $= P/$NTM EPS; \; PEG $=$ (P/E)/LT growth\\
MVE/Pretax Income: avoids tax distortions.\\
\textbf{TEV} $=$ MVE $+$ LTD $+$ Curr.LTD $-$ Cash; \; TEV/EBITDA or TEV/EBIT.\\
\textbf{Consistency}: entity-level numerator (TEV) $\leftrightarrow$ pre-debt denominator (EBIT/EBITDA); equity-level (MVE) $\leftrightarrow$ post-debt (NI, EPS, pretax).\\
\con{Special items distort: transitory vs.\ recurring? Single-peer caveat: more comps $=$ stronger conclusion. Watch units (millions vs. thousands).}

%----- SEGMENTS / NCI -----
\sect{Segments \& NCI (S9)}
Reportable segment if any: rev $\ge10\%$, profit $\ge10\%$, or assets $\ge10\%$ of all segs. Combined reported segs $\ge75\%$ of consol.\ revenue.\\
Required disclosures: P\&L per seg; often rev/assets/capex/dep; also product/geo/customer.\\
Intersegment transactions eliminated in consolidation.\\
NCI(IS) $=$ minority share of sub NI; NCI(BS) $=$ minority share of sub net assets (in equity).\\
Redeemable NCI $\to$ liabilities or temp.\ equity.\\
\con{Check: Total Equity $-$ NCI $\to$ NI to parent. NCI $=$ ``a segment the company doesn't own.'' Treat NI and equity consistently.}

%----- FORECASTING -----
\sect{Forecasting Revenue \& Expenses (S10--11)}
Revenue $=$ P $\times$ Q. Growth mean-reverts to \textbf{3--7\%}; typical 7--11\%; wide variance.\\
\textbf{Revenue recognition (5-step, FASB/IASB eff.\ Jan 2018)}:
(1) Identify contract; (2) Identify performance obligations; (3) Determine transaction price; (4) Allocate price to obligations; (5) Recognize when/as each obligation satisfied.\\
Tricky for: subscriptions, bundles, commissions, licenses, long-term contracts.\\[1pt]
Expenses: \%Sales or regression: $\text{Exp} = a + b\cdot S$. Adj for scale ($\downarrow$), lumpiness ($\uparrow$ST), sticky costs ($\uparrow$ST in downturns).\\
Dep.Exp $\approx (1/\text{Dep.Life})\times\text{Avg.Gross PPE}$; apply to existing PPE $+$ capex.\\
\con{Revenue is the engine; all forecasts follow from it. FS articulate: IS $\to$ RE $\to$ CFS. Higher fixed costs $\Rightarrow$ more margin volatility.}

%----- STOCK COMP -----
\sect{Stock Compensation (S11)}
Fair value expensed as earned (options: Black-Scholes; RSUs: price at grant). Real cost.\\
IS: salary exp. \; CFS: add back as non-cash op.\ adj.\ (like D\&A). \; BS: APIC $\uparrow$, DTA, treasury stock.\\
Often excluded from Non-GAAP; Buffett/Damodaran: real operating expense; dilutes shareholders.

%----- KEY RULES -----
\sect{Key Conceptual Rules (Quiz Focus)}
\textbf{Interpreting questions}:
\begin{itemize}[noitemsep,topsep=0pt,leftmargin=*,itemsep=0pt]
  \item State formula $\to$ apply it $\to$ acknowledge alternative interpretations $\to$ conclude
  \item ``Small accrual'' can mean demand met \textit{or} supply constrained --- say both
  \item If IS charge $=$ CFS add-back $\Rightarrow$ 100\% non-cash; state why (e.g.\ prior-period deposit)
  \item Link: accounting signal $\to$ economic logic $\to$ forward-looking implication
\end{itemize}
\textbf{F/Z model philosophy}: pre-defined recipes check analyst bias; component analysis yields same insight as DuPont drill-down.\\
\textbf{Benchmarking}: accounting diffs can be legitimate; aggressive on one $\Leftrightarrow$ conservative on another; be consistent when adjusting benchmarks; make reasonable assumptions when data is missing.\\
\textbf{Piotroski}: high score $\Rightarrow$ profitability $\uparrow$, positive surprises, stock outperformance; portfolio effect $\neq$ single stock guarantee.\\
\textbf{M-Score gray zone} ($-2.22$ to $-1.78$): inconclusive --- could signal other issues even if specific manipulation not detected.\\
\textbf{Z'' gray zone} ($4.35$--$5.85$): monitor closely; one bad year can push into distress.


%----- ACRONYMS -----
\sect{Acronym Glossary}
\textbf{Financial Statements \& Locations}\\[1pt]
{\scriptsize\begin{tabular}{@{}l@{\ }l@{\ }l@{}}
\toprule
Acr. & Meaning & Where \\
\midrule
NI   & Net Income & IS (bottom line) \\
EPS  & Earnings Per Share & IS (below NI) or calc \\
EBIT & Earnings Before Interest \& Tax & Calc: pretax $+$ int.exp. \\
EBITDA & EBIT $+$ Dep. $+$ Amort. & Calc (not on stmt) \\
COGS & Cost of Goods Sold & IS \\
GM   & Gross Margin (= Sales $-$ COGS) & IS \\
SG\&A & Selling, General \& Admin.\ Exp. & IS \\
R\&D & Research \& Development Exp. & IS \\
D\&A & Depreciation \& Amortization & CFS (op.\ add-back) / FN \\
SBC  & Stock-Based Compensation & IS (various lines) + CFS \\
OCI  & Other Comprehensive Income & Stmt of Comp.\ Inc.\ / Eq. \\
CFO  & Cash Flow from Operations & CFS (operating section) \\
CFI  & Cash Flow from Investing & CFS (investing section) \\
CFF  & Cash Flow from Financing & CFS (financing section) \\
TA   & Total Assets & BS \\
CA   & Current Assets & BS \\
CL   & Current Liabilities & BS \\
WC   & Working Capital (CA $-$ CL) & BS (calc) \\
AR   & Accounts Receivable & BS (current assets) \\
Inv  & Inventory & BS (current assets) \\
AP   & Accounts Payable & BS (current liabilities) \\
PPE  & Property, Plant \& Equipment & BS (non-current assets) \\
LTD  & Long-Term Debt & BS (long-term liabilities) \\
RE   & Retained Earnings & BS (equity section) \\
BVE  & Book Value of Equity & BS (equity section) \\
BVL  & Book Value of Liabilities & BS (total liabilities) \\
MVE  & Market Value of Equity & Calc: price $\times$ shares \\
TEV  & Total Enterprise Value & Calc: MVE$+$LTD$+$CurrLTD$-$Cash \\
NCI  & Noncontrolling Interest & IS (below NI) + BS (equity) \\
APIC & Additional Paid-In Capital & BS (equity section) \\
DTA  & Deferred Tax Asset & BS (assets) \\
LIFO & Last In, First Out (inventory) & FN (inventory note) \\
FIFO & First In, First Out (inventory) & FN (inventory note) \\
\bottomrule
\end{tabular}}\\[3pt]
\textbf{Model Variables}\\[1pt]
{\scriptsize\begin{tabular}{@{}l@{\ }l@{}}
\toprule
Acr. & Meaning \\
\midrule
DSRI  & Days Sales in Receivables Index (M) \\
GMI   & Gross Margin Index (M) \\
AQI   & Asset Quality Index (M) \\
SGI   & Sales Growth Index (M) \\
DEPI  & Depreciation Index (M) \\
SGAI  & SG\&A Expenses Index (M) \\
LVGI  & Leverage Index (M) \\
TATA  & Total Accruals to Total Assets (M) \\
NOA   & Net Operating Assets (F-Score) \\
rsst\_acc & $\Delta$Non-cash NOA / Avg.TA (F) \\
ch\_rec  & $\Delta$Receivables / Avg.TA (F) \\
ch\_inv  & $\Delta$Inventory / Avg.TA (F) \\
soft\_assets & (TA$-$PPE$-$Cash)/TA (F) \\
ch\_cs   & \% $\Delta$ Cash Sales (F) \\
ch\_roa  & $\Delta$ Return on Assets (F) \\
iss      & Debt or equity issuance (F) \\
\bottomrule
\end{tabular}}\\[3pt]
\textbf{Standards \& Regulatory}\\[1pt]
{\scriptsize\begin{tabular}{@{}l@{\ }l@{}}
\toprule
Acr. & Meaning \\
\midrule
GAAP  & Generally Accepted Accounting Principles \\
IFRS  & International Financial Reporting Standards \\
FASB  & Financial Accounting Standards Board (US) \\
IASB  & International Accounting Standards Board \\
SEC   & Securities \& Exchange Commission \\
10-K  & Annual report (IS, BS, CFS, FN in Item 8) \\
MD\&A & Management Discussion \& Analysis (Item 7) \\
AAER  & SEC Accounting \& Auditing Enforcement Release \\
ROE   & Return on Equity (calc from IS + BS) \\
ROA   & Return on Assets (calc from IS + BS) \\
P/E   & Price-to-Earnings (mkt price / EPS) \\
PEG   & Price/Earnings-to-Growth \\
\bottomrule
\end{tabular}}

\end{multicols}
\end{document}
